\documentclass[10pt,a4paper]{article}
\usepackage{blindtext}
\usepackage{subcaption}
\usepackage{graphicx}
\usepackage{tikz}
\usepackage{amssymb}
\usepackage{caption}
\usepackage{amsmath}
\usepackage{circuitikz}
\usepackage{hyperref}
\usepackage{amssymb}
\usepackage{amsmath}
\usepackage[spanish,activeacute,es-tabla]{babel}
\usepackage[utf8]{inputenc}
\usepackage{ifthen}
\usepackage{listings}
\usepackage{dsfont}
\usepackage{subcaption}
\usepackage{amsmath}
\usepackage[strict]{changepage}
\usepackage[top=1cm,bottom=2cm,left=1cm,right=1cm]{geometry}%
\usepackage{color}%
\newcommand{\tocarEspacios}{%
	\addtolength{\leftskip}{3em}%
	\setlength{\parindent}{0em}%
}

% Especificacion de procs

\newcommand{\In}{\textsf{in }}
\newcommand{\Out}{\textsf{out }}
\newcommand{\Inout}{\textsf{inout }}

\newcommand{\encabezadoDeProc}[4]{%
	% Ponemos la palabrita problema en tt
	%  \noindent%
	{\normalfont\bfseries\ttfamily proc}%
	% Ponemos el nombre del problema
	\ %
	{\normalfont\ttfamily #2}%
	\
	% Ponemos los parametros
	(#3)%
	\ifthenelse{\equal{#4}{}}{}{%
		% Por ultimo, va el tipo del resultado
		\ : #4}
}

\newenvironment{proc}[4][res]{%
	
	% El parametro 1 (opcional) es el nombre del resultado
	% El parametro 2 es el nombre del problema
	% El parametro 3 son los parametros
	% El parametro 4 es el tipo del resultado
	% Preambulo del ambiente problema
	% Tenemos que definir los comandos requiere, asegura, modifica y aux
	\newcommand{\requiere}[2][]{%
		{\normalfont\bfseries\ttfamily requiere}%
		\ifthenelse{\equal{##1}{}}{}{\ {\normalfont\ttfamily ##1} :}\ %
		\{\ensuremath{##2}\}%
		{\normalfont\bfseries\,\par}%
	}
	\newcommand{\asegura}[2][]{%
		{\normalfont\bfseries\ttfamily asegura}%
		\ifthenelse{\equal{##1}{}}{}{\ {\normalfont\ttfamily ##1} :}\
		\{\ensuremath{##2}\}%
		{\normalfont\bfseries\,\par}%
	}
	\renewcommand{\aux}[4]{%
		{\normalfont\bfseries\ttfamily aux\ }%
		{\normalfont\ttfamily ##1}%
		\ifthenelse{\equal{##2}{}}{}{\ (##2)}\ : ##3\, = \ensuremath{##4}%
		{\normalfont\bfseries\,;\par}%
	}
	\renewcommand{\pred}[3]{%
		{\normalfont\bfseries\ttfamily pred }%
		{\normalfont\ttfamily ##1}%
		\ifthenelse{\equal{##2}{}}{}{\ (##2) }%
		\{%
		\begin{adjustwidth}{+5em}{}
			\ensuremath{##3}
		\end{adjustwidth}
		\}%
		{\normalfont\bfseries\,\par}%
	}
	
	\newcommand{\res}{#1}
	\vspace{1ex}
	\noindent
	\encabezadoDeProc{#1}{#2}{#3}{#4}
	% Abrimos la llave
	\par%
	\tocarEspacios
}
{
	% Cerramos la llave
	\vspace{1ex}
}

\newcommand{\aux}[4]{%
	{\normalfont\bfseries\ttfamily\noindent aux\ }%
	{\normalfont\ttfamily #1}%
	\ifthenelse{\equal{#2}{}}{}{\ (#2)}\ : #3\, = \ensuremath{#4}%
	{\normalfont\bfseries\,;\par}%
}

\newcommand{\pred}[3]{%
	{\normalfont\bfseries\ttfamily\noindent pred }%
	{\normalfont\ttfamily #1}%
	\ifthenelse{\equal{#2}{}}{}{\ (#2) }%
	\{%
	\begin{adjustwidth}{+2em}{}
		\ensuremath{#3}
	\end{adjustwidth}
	\}%
	{\normalfont\bfseries\,\par}%
}

% Tipos

\newcommand{\nat}{\ensuremath{\mathds{N}}}
\newcommand{\ent}{\ensuremath{\mathds{Z}}}
\newcommand{\float}{\ensuremath{\mathds{R}}}
\newcommand{\bool}{\ensuremath{\mathsf{Bool}}}
\newcommand{\cha}{\ensuremath{\mathsf{Char}}}
\newcommand{\str}{\ensuremath{\mathsf{String}}}

% Logica

\newcommand{\True}{\ensuremath{\mathrm{true}}}
\newcommand{\False}{\ensuremath{\mathrm{false}}}
\newcommand{\Then}{\ensuremath{\rightarrow}}
\newcommand{\Iff}{\ensuremath{\leftrightarrow}}
\newcommand{\implica}{\ensuremath{\longrightarrow}}
\newcommand{\IfThenElse}[3]{\ensuremath{\mathsf{if}\ #1\ \mathsf{then}\ #2\ \mathsf{else}\ #3\ \mathsf{fi}}}
\newcommand{\yLuego}{\land _L}
\newcommand{\oLuego}{\lor _L}
\newcommand{\implicaLuego}{\implica _L}

\newcommand{\cuantificador}[5]{%
	\ensuremath{(#2 #3: #4)\ (%
		\ifthenelse{\equal{#1}{unalinea}}{
			#5
		}{
			$ % exiting math mode
			\begin{adjustwidth}{+2em}{}
				$#5$%
			\end{adjustwidth}%
			$ % entering math mode
		}
		)}
}

\newcommand{\existe}[4][]{%
	\cuantificador{#1}{\exists}{#2}{#3}{#4}
}
\newcommand{\paraTodo}[4][]{%
	\cuantificador{#1}{\forall}{#2}{#3}{#4}
}

%listas

\newcommand{\TLista}[1]{\ensuremath{seq \langle #1\rangle}}
\newcommand{\lvacia}{\ensuremath{[\ ]}}
\newcommand{\lv}{\ensuremath{[\ ]}}
\newcommand{\longitud}[1]{\ensuremath{|#1|}}
\newcommand{\cons}[1]{\ensuremath{\mathsf{addFirst}}(#1)}
\newcommand{\indice}[1]{\ensuremath{\mathsf{indice}}(#1)}
\newcommand{\conc}[1]{\ensuremath{\mathsf{concat}}(#1)}
\newcommand{\cab}[1]{\ensuremath{\mathsf{head}}(#1)}
\newcommand{\cola}[1]{\ensuremath{\mathsf{tail}}(#1)}
\newcommand{\sub}[1]{\ensuremath{\mathsf{subseq}}(#1)}
\newcommand{\en}[1]{\ensuremath{\mathsf{en}}(#1)}
\newcommand{\cuenta}[2]{\mathsf{cuenta}\ensuremath{(#1, #2)}}
\newcommand{\suma}[1]{\mathsf{suma}(#1)}
\newcommand{\twodots}{\ensuremath{\mathrm{..}}}
\newcommand{\masmas}{\ensuremath{++}}
\newcommand{\matriz}[1]{\TLista{\TLista{#1}}}
\newcommand{\seqchar}{\TLista{\cha}}

\renewcommand{\lstlistingname}{Código}
\lstset{% general command to set parameter(s)
	language=Java,
	morekeywords={endif, endwhile, skip},
	basewidth={0.47em,0.40em},
	columns=fixed, fontadjust, resetmargins, xrightmargin=5pt, xleftmargin=15pt,
	flexiblecolumns=false, tabsize=4, breaklines, breakatwhitespace=false, extendedchars=true,
	numbers=left, numberstyle=\tiny, stepnumber=1, numbersep=9pt,
	frame=l, framesep=3pt,
	captionpos=b,
}

\newcommand{\notimplies}{\;\not\!\!\!\implies}
\title{Algoritmos y Estructuras de Datos II}
\author{Tomás Agustín Hernández}
\date{}

\begin{document}
\maketitle

\begin{figure}[b]
    \centering
    \begin{tikzpicture}[remember picture,overlay]
        \node[anchor=south east, inner sep=0pt, xshift=-1cm, yshift=2cm] at (current page.south east) {
            \begin{minipage}[b]{0.5\textwidth}
                \includegraphics[width=\linewidth]{logo_uba.jpg}
                \label{fig:bottom}
            \end{minipage}
        };
    \end{tikzpicture}
\end{figure}

\newpage
\section{Especificación}
\subsection*{Consideraciones importantes / Reminders}
\begin{itemize}
    \item Utilizar operadores luego: Si estoy en LPO (Lógica de Primer Orden) utilizar los operadores luego si vemos que hay una posible indefinición como una división, o ingresar a una lista a un índice. Recordar que el para todo y un existe, aunque esté acotado por un rango, los cuantificadores predican IGUAL para todos los valores. Entonces, aunque diga que x es positivo, también probará dividir inclusive por 0 y estallará.
    \item Recordar las condiciones bidireccionales
        \begin{itemize}
            \item Si por algún motivo tengo que armar una “lista”, como, por ejemplo, los divisores de un número x tengo que indicar que, si el número divide a x, entonces ese número está en res, pero además todos los valores que están en res DIVIDEN a x. Es una condición bidireccional. 
            \item Otro ejemplo puede ser que tenga que considerar el máximo de una lista, si todos los valores y que están en la lista son menores que res entonces significa que res también pertenece a esa lista original.
        \end{itemize}
    \item Recordar el significado de los cuantificadores con dos variables al mismo tiempo: En la lógica se ejecutan todos de uno a la vez. Es decir, si tengo que poner un para todo adentro de un para todo entonces hago un para todo solo con dos variables y listo.
    \item Recordar que cuando en un procedimiento llamo a un predicado y ese predicado devuelve algo de un para todo, existe (básicamente un valor de verdad) tengo que castear ese valor en el procedimiento porque son dos mundos distintos.
    Ej: asegura: { res = True \(\iff\) predicado}
    \item Los predicados y funciones auxiliares no describen problemas. Son herramientas sintácticas para descomponer predicados.
    \begin{itemize}
        \item Los procedimientos pueden llamar a funciones auxiliares o predicados. Un procedimiento no puede llamar a otro procedimiento.
        \item Los predicados pueden llamar a predicados o auxiliares. 
        \item Las auxiliares solo pueden llamar auxiliares.
    \end{itemize}
    \item No usamos nunca \(==\) en especificación, usamos siempre \(=\) y estamos comparando, no asignando.
    \item No existe el guardar o asignar en el mundo de la lógica. No puedo guardar en una lista en un índice específico porque si un valor. Para esto solemos usar que x valor pertenecerá a esta lista, por ejemplo.
    \item Si tengo un algoritmo que cumple una funcionalidad específica con un require más débil, puedo poner el require más restrictivo y va a funcionar igual pero NO al revés.
\end{itemize}
\subsection*{Fórmulas compuestas}
Decimos que una fórmula es compuesta a una fórmula que tiene más de una operación y esa operación necesita realizarse antes de conocer su valor.
\begin{itemize}
    \item \((p \land q) \lor m\)
    \item \(((p \land q) \lor m) \implies n\)
\end{itemize}
\subsection*{Fórmula atómica}
Decimos que una fórmula es atómica si se puede inferir su valor con una, o ninguna operación. Es irreducible.
\begin{itemize}
    \item p
    \item \(p \land q\)
\end{itemize}
\subsection*{Fórmulas bien definidas}
Decimos que una fórmula está bien definida cuando el orden que hay que hacer las operaciones es clara. Es decir, cuando cada operación toma dos variables proposicionales, y al realizar la operación termina siendo una fórmula atómica.
\begin{itemize}
    \item \(p \land q \lor r\) está mal formada. No se especifica si primero se realiza el \(\land\) o el \(\lor\).
    \item \((p \land q) \lor r\) está bien formada.
    \item \(p \land q \land r \land m\) está bien formada porque son todas conjunciones.
    \item \(p \lor q \lor r \lor m\) está bien formada porque son todas disyunciones.
\end{itemize}
\subsection*{Cuantificadores}
\begin{itemize}
    \item Para todo: \(\forall\)
    \begin{itemize}
        \item \(Garantiza \ la \ \text{conjunción}: p(1) \land p(2) \land p(3) \dots \land p(m) \). Todos los casos deben ser \True \ para que el cuantificador sea \True.
        \item Se acompaña por un \(\implica\) a la hora de predicar sobre los elementos.
        \item \((\forall i: \ent)(0 \le i < \longitud{s} \implicaLuego \ s[i] \ mod \ 2 = 0)\). Todos los elementos de la lista son divisibles por 2.
    \end{itemize}
    \item Existe: \(\exists\)
    \begin{itemize}
        \item \(Garantiza \ la \ \text{disyunción}: p(1) \lor p(2) \lor p(3) \dots \lor p(m) \). Con un caso \True \ el cuantificador es \True.
        \item Se acompaña por un \(\land\) a la hora de predicar sobre los elementos.
        \item \((\exists i: \ent)(0 \le i < \longitud{s} \yLuego s[i] \ge 0)\). Existe algún elemento en la lista que es mayor o igual a 0.
    \end{itemize}
\end{itemize}

\subsection*{Equivalencias entre fórmulas}
Decimos que dos fórmulas son equivalentes \(\iff\) los valores de la tabla de verdad al aplicar la operación arroja el mismo resultado.

\subsection*{Valuaciones}
Las valuaciones surgen en base a la tabla de verdad. Las valuaciones serian darle valor a las variables proposicionales y ver el resultado de la operación. Solo hacen referencias a fórmulas atómicas.

\subsection*{Tautologias, contradicciones y contingencias}
\begin{itemize}
    \item Una fórmula es tautología \(\iff\) el resultado de la operación en cada fila arroja siempre V.
    \item Una fórmula es contradicción \(\iff\) el resultado de la operación en cada fila arroja siempre F.
    \item  Una fórmula es contradicción \(\iff\) el resultado de la operación en cada fila arroja siempre V y F.
\end{itemize}

\subsection*{Relaciones de fuerza entre fórmulas}
Decimos que una fórmula es más fuerte que la otra \(\iff\) una fórmula es más restrictiva que la otra, o está incluida en la otra. \\
En el mundo de la lógica, decimos que A es más fuerte que B \(\iff A \implies B\)

\begin{itemize}
    \item Si (\(A \implies B\)) y (\(B \implies A\)) son tautologías, entonces A y B son equivalentes.
    \item Si (\(A \implies B\)) es tautología y (\(B \notimplies A\)) no es tautología, entonces decimos que A es más fuerte que B.
    \item Si (\(A \notimplies B\)) y (\(B \notimplies A\)) son contigencias, entonces no existe relación de fuerza entre A y B.
\end{itemize}

Algunos ejemplos:

\begin{itemize}
    \item \(\longitud{s} = 0 \implies \longitud{s} \ge 0\). En este caso vemos que \(\longitud{s} = 0\) es más fuerte que \(\longitud{s} \ge 0\) pues \(\longitud{s} = 0\) está incluido en \(\longitud{s} \ge 0\). Por lo tanto, \(A \implies B\)
    \item \(\longitud{s} = 0 \implies \longitud{s} \ge 3\). En este caso vemos que \(\longitud{s} = 0\) no es más fuerte que \(\longitud{s} \ge 3\) pues \(\longitud{s} = 0\) no está incluido en \(\longitud{s} \ge 3\). Por lo tanto, \(A \notimplies B\)
    \item \(2 \le i < \longitud{s} \implies 1 \le i < \longitud{s}\). En este caso A \(\implies\) B, pues i = 2 está incluido en el rango de B. Por lo tanto, \(A \implies B\)
    \item \(0 \le i < \longitud{s} \implies 1 \le i < \longitud{s}\). En este caso A \(\notimplies\) B, pues el 0 de A no es parte de B. Por lo tanto, \(A \notimplies B\)
\end{itemize}

\subsection*{Lógica trivaluada}
También llamada lógica secuencial porque se procesa de izquierda a derecha; Nos introduce los conceptos de \(\yLuego \ \oLuego \ \implicaLuego\) y el valor de indefinido \(\bot\).

Se termina de evaluar una expresión cuando se puede deducir el valor de verdad. \\


Considere \(x = \True \land y = \bot \land z = \False \)
\begin{itemize}
    \item \(x \oLuego y\): Como el \(\oLuego\) necesita uno solo para ser verdadero, entonces como x ya es \(\True\) entonces toda la fórmula es verdadera.
    \item \(x \yLuego y\): Como el \(\yLuego\) necesita que ambas variables sean verdaderas, evalúa indefinido y el programa estalla.
    \item \(\neg x \implicaLuego y\): Como el \(\implicaLuego\) solo es falso si el antecedente es  \(\True\) y el consecuente \(\False\), como en este caso el antecedente ya es falso, toda la implicación es verdadera.
    \item \((x \land z) \yLuego y\): Como el \(\yLuego\) necesita que ambas fórmulas sean \(\True\), en este caso, como \((x \land z)\) es falso, entonces ya toda la fórmula es falsa. Nótese que el \(\land\) de la condición interna no contiene el luego porque jamás se indefinirá.
    \item \((\forall i: \ent)(0 \le i < \longitud{s} \implicaLuego s[i] \ge 0)\) Nótese que aquí usamos un \(\implicaLuego\) porque podría ser que la lista esté indefinida o no exista el valor en s[i]
\end{itemize}



\end{document}